\section{Бинарное возведение в степень по модулю}
Вычисляет $a^e\bmod m$ за $O(\log(e+1))$ времени и $O(1)$ памяти, возвращая число из $[0,m)$.
Основание --- любое значение \texttt{int}, $0\le e\le\texttt{LLONG\_MAX}$, модуль типа \texttt{int}: $1\le m\le2\cdot10^9$; произведения вычисляются в \texttt{ll} и меньше $4\cdot10^{18}$.
При $e=0$ результат равен $1\bmod m$, в том числе для $a=0$.
\begin{lstlisting}[name=math_power]
int modularPower(int base, ll exponent, int modulus) {
    assert(exponent >= 0 && modulus > 0 && modulus <= 2000000000);
    base %= modulus;
    if (base < 0) base += modulus;
    int result = 1 % modulus;
    while (exponent > 0) {
        if (exponent & 1)
            result = (ll)result * base % modulus;
        base = (ll)base * base % modulus;
        exponent >>= 1;
    }
    return result;
}
\end{lstlisting}
% \subsection*{Для чего нужен}
% Вычисляет $a^e\bmod m$ за $O(\log(e+1))$ времени и $O(1)$ памяти.
% Основание --- любое значение \texttt{int}, $0\le e\le\texttt{LLONG\_MAX}$,
% модуль типа \texttt{int}: $1\le m\le2\cdot10^9$. Результат лежит в $[0,m)$.
% При $e=0$ возвращает $1\bmod m$, в том числе для $a=0$.
% Произведения вычисляются в \texttt{ll} и меньше $4\cdot10^{18}$.
% \subsection*{Реализация}
% Материал: \url{https://ru.algorithmica.org/cs/algebra/binpow/}.
%
% \subsection*{Возможные улучшения и дополнения}
% Для отрицательного показателя сначала найдите обратный элемент;
% он существует только при взаимной простоте основания и модуля.
% Сокращать показатель по $\varphi(m)$ можно при $\gcd(a,m)=1$.
% При большом числе операций с одним модулем возможна редукция Монтгомери.

\section{Расширенный алгоритм Евклида и обратный элемент}
Расширенный алгоритм Евклида возвращает тройку $(g,x,y)$ с $g=\gcd(a,b)$ и $ax+by=g$ для входов \texttt{int}, $0\le a,b\le\texttt{INT\_MAX}$; для $(0,0)$ результат $(0,1,0)$.
Обратный элемент для любого \texttt{int} по модулю $2\le m\le\texttt{INT\_MAX}$ возвращается как \texttt{optional<int>}: пустой при $\gcd(value,m)\ne1$, иначе число из $[0,m)$.
Коэффициенты и промежуточные произведения вычисляются в \texttt{ll}: модули коэффициентов не превосходят \texttt{INT\_MAX}, произведение частного и коэффициента --- $\texttt{INT\_MAX}^2$, вычитание также безопасно.
Время $O(\log(\max(a,b)+2))$, память $O(1)$.
\begin{lstlisting}[name=math_euclid]
tuple<int, ll, ll> extendedGcd(int first, int second) {
    assert(first >= 0 && second >= 0);
    ll firstX = 1, firstY = 0, secondX = 0, secondY = 1;
    while (second != 0) {
        int quotient = first / second;
        int remainder = first % second;
        first = second;
        second = remainder;
        ll nextX = firstX - quotient * secondX;
        ll nextY = firstY - quotient * secondY;
        firstX = secondX;
        firstY = secondY;
        secondX = nextX;
        secondY = nextY;
    }
    return {first, firstX, firstY};
}

optional<int> modularInverse(int value, int modulus) {
    assert(modulus >= 2);
    value %= modulus;
    if (value < 0) value += modulus;
    auto [divisor, coefficient, unused] = extendedGcd(value, modulus);
    if (divisor != 1) return nullopt;
    coefficient %= modulus;
    if (coefficient < 0) coefficient += modulus;
    return coefficient;
}
\end{lstlisting}
% \subsection*{Для чего нужен}
% Находит $g=\gcd(a,b)$ и коэффициенты $ax+by=g$ для
% $0\le a,b\le\texttt{INT\_MAX}$, входы имеют тип \texttt{int}.
% Для $(0,0)$ возвращает $(0,1,0)$; результат --- тройка $(g,x,y)$.
% Время $O(\log(\max(a,b)+2))$, память $O(1)$.
% Обратный элемент ищется для произвольного \texttt{int} по модулю
% $2\le m\le\texttt{INT\_MAX}$. Возвращает \texttt{optional<int>}:
% пустой при $\gcd(value,m)\ne1$, иначе содержит число из $[0,m)$.
% Все коэффициенты и их промежуточные произведения вычисляются в \texttt{ll}.
% При этих входах модули коэффициентов не превосходят \texttt{INT\_MAX},
% а произведение частного и коэффициента --- $\texttt{INT\_MAX}^2$;
% вычитание также безопасно в \texttt{ll}.
% \subsection*{Реализация}
% Материал: \url{https://ru.algorithmica.org/cs/modular/extended-euclid/},
% \url{https://ru.algorithmica.org/cs/modular/reciprocal/}.
%
% \subsection*{Возможные улучшения и дополнения}
% Для $ax+by=c$ необходимо $g\mid c$; при $g>0$ умножьте коэффициенты
% на $c/g$, контролируя переполнение. Все решения:
% $(x+k\,b/g,y-k\,a/g)$ после такого масштабирования.
% Алгоритм также используется в китайской теореме об остатках.
% Для простого модуля обратный ненулевого остатка равен $a^{m-2}\bmod m$.

\section{Линейное решето с минимальным простым делителем}
Строит массив минимальных простых делителей чисел от $2$ до входного $n$; для $0$ и $1$ записан ноль, при $x\ge2$ число простое тогда и только тогда, когда \texttt{smallest[x] == x}.
Требуется $0\le n<\texttt{INT\_MAX}$ и достаточно памяти; время и память $O(n)$.
Для факторизации $2\le x\le n$ последовательно берут \texttt{smallest[x]} и делят $x$ на него за $O(\log x)$.
\begin{lstlisting}[name=math_sieve]
vector<int> smallestPrimeFactors(int limit) {
    assert(limit >= 0 && limit < INT_MAX);
    vector<int> smallest(limit + 1), primes;
    for (int number = 2; number <= limit; ++number) {
        if (smallest[number] == 0) {
            smallest[number] = number;
            primes.push_back(number);
        }
        for (int prime : primes) {
            if (prime > smallest[number] || prime > limit / number) break;
            smallest[number * prime] = prime;
        }
    }
    return smallest;
}
\end{lstlisting}
% \subsection*{Для чего нужен}
% Строит массив минимальных простых делителей всех чисел от $2$ до $n$.
% Время $O(n)$, память $O(n)$; требуется $0\le n<\texttt{INT\_MAX}$
% и достаточно памяти. Для $0$ и $1$ записан ноль.
% Число $x\ge2$ простое тогда и только тогда, когда
% \texttt{smallest[x] == x}. Для факторизации $2\le x\le n$ последовательно
% берите \texttt{smallest[x]} и делите $x$ на него; время $O(\log x)$.
% \subsection*{Реализация}
% Материал: \url{https://ru.algorithmica.org/cs/factorization/eratosthenes/}.
%
% \subsection*{Возможные улучшения и дополнения}
% В том же проходе можно вычислять функцию Эйлера и функцию Мёбиуса.
% Если нужны только признаки простоты, битовое решето Эратосфена экономит
% память. Для большого интервала $[L,R]$ используйте сегментированное решето
% с простыми до $\sqrt R$; хранить делители всех чисел до $R$ не потребуется.

\section{Метод Гаусса для решения линейной системы}
Решает вещественную систему $Ax=b$ по расширенной матрице из $n$ строк по $m+1$ чисел с $b$ в последнем столбце; возвращает статус $0$ (нет решения, вектор пуст), $1$ (одно) или $2$ (бесконечно много, возвращено одно со свободными переменными, равными нулю).
Допустимы $n=0$, $m=0$; размеры меньше \texttt{INT\_MAX}, все числа конечны, промежуточные вычисления не должны переполнять \texttt{db}.
Порог \texttt{epsilon} конечный и положительный, сравнения ведущих элементов и невязки абсолютные; данные следует масштабировать, для плохо обусловленных систем статус и решение могут быть неточными.
Время $O(nm\min(n,m)+nm)$, память $O(n(m+1)+m)$ с копией входа.
\begin{lstlisting}[name=math_gauss]
pair<int, vector<db>> gauss(vector<vector<db>> matrix, int variables,
                            db epsilon = 1e-9) {
    assert(variables >= 0 && variables < INT_MAX);
    assert(isfinite(epsilon) && epsilon > 0 && matrix.size() < INT_MAX);
    int rows = matrix.size(), rank = 0;
    for (const auto& row : matrix) {
        assert(row.size() == static_cast<unsigned>(variables + 1));
        for (db value : row) assert(isfinite(value));
    }
    auto original = matrix;
    vector<int> pivotRow(variables, -1);
    for (int column = 0; column < variables && rank < rows; ++column) {
        int selected = rank;
        for (int row = rank + 1; row < rows; ++row)
            if (abs(matrix[row][column]) > abs(matrix[selected][column]))
                selected = row;
        if (abs(matrix[selected][column]) <= epsilon) continue;
        swap(matrix[selected], matrix[rank]);
        pivotRow[column] = rank;
        for (int entry = variables; entry >= column; --entry)
            matrix[rank][entry] /= matrix[rank][column];
        for (int row = 0; row < rows; ++row) {
            if (row == rank) continue;
            for (int entry = variables; entry >= column; --entry)
                matrix[row][entry] -= matrix[rank][entry] * matrix[row][column];
        }
        ++rank;
    }
    vector<db> solution(variables);
    for (int column = 0; column < variables; ++column)
        if (pivotRow[column] != -1)
            solution[column] = matrix[pivotRow[column]][variables];
    for (const auto& row : original) {
        db residual = -row[variables];
        for (int column = 0; column < variables; ++column)
            residual += row[column] * solution[column];
        if (abs(residual) > epsilon) return {0, {}};
    }
    return {rank == variables ? 1 : 2, solution};
}
\end{lstlisting}
% \subsection*{Для чего нужен}
% Решает $Ax=b$ с вещественными коэффициентами. Вход --- расширенная
% матрица из $n$ строк по $m+1$ чисел, последний столбец содержит $b$.
% Статусы: $0$ --- нет решения (вектор пуст), $1$ --- одно,
% $2$ --- бесконечно много (возвращено одно, свободные переменные равны нулю).
% Допустимы $n=0$ и $m=0$; размеры меньше \texttt{INT\_MAX}.
% Время $O(nm\min(n,m)+nm)$, память $O(n(m+1)+m)$ с копией входа.
% Все числа конечны; промежуточные вычисления не должны переполнять \texttt{db}.
% Порог \texttt{epsilon} положительный, абсолютный: малые ведущие элементы
% считаются нулём. Масштабируйте данные; для плохо обусловленных систем
% статус и решение могут быть неточными. Проверка невязки также абсолютная.
% \subsection*{Реализация}
% Материал: \url{https://ru.algorithmica.org/cs/algebra/gauss/};
% реализация адаптирована для прямоугольных и вырожденных систем.
% Столбцы обрабатываются справа налево, чтобы ведущий элемент
% и множитель строки сохранялись до последнего вычитания.
%
% \subsection*{Возможные улучшения и дополнения}
% Для описания всех решений постройте базис ядра по свободным переменным.
% Для точной арифметики по простому модулю замените деление обратным элементом;
% над составным модулем этот подход напрямую не работает.
% Относительный порог, \texttt{long double} и выбор ведущего элемента
% по всей подматрице могут повысить устойчивость, но не гарантируют её.

\section{Выпуклая оболочка целочисленных точек}
Монотонная цепь строит минимальный выпуклый многоугольник, содержащий входные точки \texttt{pair<int, int>}; результат идёт против часовой стрелки от лексикографически минимальной точки без её повторения, дубликаты и промежуточные точки сторон удалены.
Требуется $|x|,|y|\le10^9$ и число точек меньше \texttt{INT\_MAX}; приведение к \texttt{ll} делается до вычитания, модули разностей не больше $2\cdot10^9$, произведений --- $4\cdot10^{18}$, их разности --- $8\cdot10^{18}$, что помещается в \texttt{ll}.
Пустой вход даёт пустой ответ, одна уникальная точка возвращается сама, коллинеарный набор --- два крайних конца в лексикографическом порядке.
Время $O(n\log(n+1))$, память $O(n)$.
\begin{lstlisting}[name=geometry_hull]
vector<pair<int, int>> convexHull(vector<pair<int, int>> points) {
    assert(points.size() < INT_MAX);
    for (auto [x, y] : points) {
        assert(-1000000000 <= x && x <= 1000000000);
        assert(-1000000000 <= y && y <= 1000000000);
    }
    sort(points.begin(), points.end());
    points.erase(unique(points.begin(), points.end()), points.end());
    if (points.size() <= 1) return points;
    auto cross = [](const auto& first, const auto& second, const auto& third) {
        ll dx = (ll)second.first - first.first;
        ll dy = (ll)second.second - first.second;
        return dx * ((ll)third.second - first.second)
             - dy * ((ll)third.first - first.first);
    };
    vector<pair<int, int>> lower, upper;
    for (const auto& point : points) {
        while (lower.size() >= 2 &&
               cross(lower[lower.size() - 2], lower.back(), point) <= 0)
            lower.pop_back();
        lower.push_back(point);
    }
    for (auto iterator = points.rbegin(); iterator != points.rend(); ++iterator) {
        while (upper.size() >= 2 &&
               cross(upper[upper.size() - 2], upper.back(), *iterator) <= 0)
            upper.pop_back();
        upper.push_back(*iterator);
    }
    lower.pop_back();
    upper.pop_back();
    lower.insert(lower.end(), upper.begin(), upper.end());
    return lower;
}
\end{lstlisting}
% \subsection*{Для чего нужен}
% Алгоритм монотонной цепи строит минимальный выпуклый многоугольник,
% содержащий все точки. Точка --- \texttt{pair<int, int>} с координатами $(x,y)$;
% $|x|,|y|\le10^9$, число точек меньше \texttt{INT\_MAX}.
% Приведение к \texttt{ll} выполняется до вычитания координат.
% Модули разностей не больше $2\cdot10^9$, произведений --- $4\cdot10^{18}$,
% их разности --- $8\cdot10^{18}$, что помещается в \texttt{ll}.
% Время $O(n\log(n+1))$, память $O(n)$.
% Результат идёт против часовой стрелки от лексикографически минимальной точки,
% без повторения первой в конце. Дубликаты и промежуточные точки сторон удалены.
% Пустой вход даёт пустой ответ; для одной уникальной точки возвращается она,
% для коллинеарных точек --- два крайних конца в лексикографическом порядке.
% \subsection*{Реализация}
% Материал о построении оболочки стеком:
% \url{https://ru.algorithmica.org/cs/convex-hulls/graham/}.
% Здесь используется монотонная цепь, а не полярная сортировка из статьи.
%
% \subsection*{Возможные улучшения и дополнения}
% Чтобы сохранить все точки на границе, замените условие поворота
% $\le0$ на $<0$ и отдельно обработайте полностью коллинеарный набор,
% иначе точки будут повторяться. По оболочке можно искать диаметр
% вращающимися калиперами за $O(h)$, где $h$ --- число её вершин.
% При расширении диапазона координат заново проверьте границы произведений:
% арифметика \texttt{ll} может переполниться.

